\par\Needspace{6\baselineskip}\medskip
\begin{center}\hypertarget{native-172-article125}{}\pdfbookmark[2]{125.}{article-125}125.\end{center}
With a slight modification of the preceding method we obtain a \emph{second} one. Starting from the distances at two observations, we determine all the elements just as before. For the third observation, the calculation from these elements goes only as far as the heliocentric position in the orbit, leaving the geocentric position uncomputed. We also derive the same heliocentric position from the observed geocentric position and the orientation of the orbital plane by solving the problem treated in Articles 74 and 75. These two determinations, which differ unless the true values of $x,y$ happen to have been assumed, supply $X,Y$. We take as $X$ the difference between the two values of the longitude in the orbit and as $Y$ the \SourceLeaf{173}difference between the two values of the radius vector, or rather of its logarithm. The same cautions apply to this method as were noted in the preceding article. We must add a further one. The heliocentric position in the orbit cannot be derived from the geocentric position whenever the position of the Earth coincides with either node of the orbit. This method is therefore inapplicable in that case. We should also avoid it when the position of the Earth is very close to either node, because the assumption that small variations of $x,y$ give rise to proportional variations of $X,Y$ would be too inaccurate, for a reason similar to the one noted in the preceding article. Here again we may seek a remedy by exchanging the middle position with one of the two extreme positions whose corresponding position of the Earth lies farther from the nodes, unless the Earth happens to have been in the vicinity of the nodes at all three observations.
